\chapter*{investigaci\'on GRUPAL}
\label{cap:cap1}
\vspace{-10ex}

%\begin{flushright}
%\linea{1.2}
%\end{flushright}
%\thispagestyle{empty}
\raggedright

%\section*{Aspectos Generales:} 

Este proyecto tiene como objetivo que el estudiantado relacione métodos matemáticos a temas de vanguardia en la física y la ingeniería. A través de la investigación y el trabajo colaborativo, se espera que los equipos puedan:

\begin{itemize}
    \item \textbf{\textcolor{c1}{profundizar en el conocimiento de un área específica}}: cada equipo seleccionará un tema relevante en la física e ingeniería modernas, investigará las bases teóricas y explorará sus aplicaciones.
    
    \item \textbf{\textcolor{c1}{aplicar métodos matemáticos avanzados}}: el proyecto está diseñado para que pongan en práctica métodos matemáticos, aplicándolos a la solución de problemas relacionados con su tema de investigación.
    
    \item \textbf{\textcolor{c1}{desarrollar habilidades de comunicación técnica}}: además de la investigación, los equipos deberán redactar un informe claro y coherente que incluya una introducción, objetivos y aplicaciones, y presentarán públicamente sus resultados. 
    
    \item \textbf{\textcolor{c1}{fomentar el trabajo en equipo y la responsabilidad compartida}}: al trabajar en equipo, aprenderán a dividir tareas, gestionar tiempos, y asumir la responsabilidad por la calidad del trabajo grupal.
\end{itemize}

Los equipos deber\'an estar conformados por un \textbf{\textcolor{c1}{m\'aximo de tres (3) personas}}. El trabajo debe ser \textbf{\textcolor{c1}{original}} y cumplir con las normas acad\'emicas del Instituto Tecnol\'ogico. \\[.5cm]



Este proyecto está dividido en \textbf{\textcolor{c1}{cuatro (4) etapas}}, cada una con entregables específicos y un porcentaje asignado en la evaluación final; seg\'un se detalla a continuaci\'on.

\section*{Etapa 1 (P1): Definici\'on del equipo, del tema y de los objetivos (1 \%)}

\begin{itemize}
    \item[\checkmark] Fecha de entrega: martes 18 de marzo 2025 (SEMANA 6).
    \item[\checkmark] Extensi\'on:  m\'aximo 3 p\'aginas.
    \item[\checkmark] Formato de entrega: documento en formato .pdf a través del módulo de evaluaciones en el \td.
    \item[\checkmark] Estructura: Portada$+$Objetivos$+$Referencias.
\end{itemize}

%\newpage

\textbf{Tareas a realizar:}

\begin{enumerate}
    \item \textbf{Elegir el equipo de trabajo}: formar un equipo de máximo 3 personas. 
    \item \textbf{Seleccionar un tema de investigaci\'on}: elegir alg\'un tema entre
    \begin{itemize}
        \item Machine learning y procesamiento de señales
        \item Métodos de Monte Carlo
        \item Funciones de Green
        \item Geometr\'ia diferencial
        \item Modelado de sistemas físicos complejos 
        \item Procesamiento de imágenes y visión por computadora
        \item Caos y sistemas complejos
        \item M\'etodos asint\'oticos y teor\'ia de perturbaciones
        \item Dinámica de fluidos computacional (CFD)
        \item An\'alisis ac\'ustico
        \item Mec\'anica computacional
        \item Computaci\'on cu\'antica
    \end{itemize}
\end{enumerate}


\resp{.75\textwidth}{\textit{Varios equipos pueden seleccionar el mismo tema. También pueden proponer un tema no incluido, previa aprobación del docente. }
}
{white}{comunicacion}


\begin{enumerate}
\setcounter{enumi}{2}
\item Una vez realizados los pasos \SebastianoItemA{1} y \SebastianoItemA{2} ingresar al siguiente \href{https://estudianteccr-my.sharepoint.com/:x:/g/personal/prof_glacy_estudiantec_cr/Eb3cnsJ318ZCiCbvQZV8QToB3RXEgbzVeGjASaGxJvZx7g?e=UOcGrV}{\textbf{documento 
 de Excel}} y registrar la conformaci\'on del equipo y el tema seleccionado. \textbf{Se debe ingresar con credenciales en el dominio @estudiantec.cr}.
\item \textbf{Redactar los objetivos}: definir \textbf{un objetivo general} y \textbf{tres objetivos especificos}, que  delimiten el alcance del trabajo. 
%\item \textbf{Introducci\'on}: elaborar una breve rese\~na del tema seleccionado; incluyendo una descripci\'on de los m\'etodos m\'atem\'aticos involucrados y su campo de \textbf{aplicaci\'on}.
\item \textbf{Referencias bibliogr\'aficas}:  incluir al menos \textbf{\textcolor{c1}{tres referencias bibliogr\'aficas}}, en formato APA o IEEE.

\end{enumerate}


\section*{Etapa 2 (P2): Primer avance (3 \%)}

\begin{itemize}
    \item[\checkmark] Fecha de entrega: viernes 11 de abril 2025 (SEMANA 8).
    \item[\checkmark] Extensi\'on:  m\'aximo 8 p\'aginas.
    \item[\checkmark] Formato de entrega: documento en formato .pdf a través del módulo de evaluaciones en el \td.
       
    \item[\checkmark] Estructura: 
     Elementos de la Etapa 1$+$Marco te\'orico.
\end{itemize}
%\newpage
\textbf{Tareas a realizar:}

\begin{enumerate}
    \item \textbf{Marco teórico}: para cada uno de los objetivos planteados en la \textbf{Etapa 1}, desarrollar un capítulo donde se exponga lo más relevante para alcanzarlos. Hacer \'enfasis en los m\'etodos matem\'aticos que sean relevantes.

\end{enumerate}

\section*{Etapa 3 (P3): Segundo avance (3 \%)}

\begin{itemize}
    \item[\checkmark] Fecha de entrega: viernes 16 de mayo 2025 (SEMANA 12).
    \item[\checkmark] Extensi\'on:  m\'aximo 12 p\'aginas.
    \item[\checkmark] Formato de entrega: documento en formato .pdf a través del módulo de evaluaciones en el \td.
       
    \item[\checkmark] Estructura: 
     Elementos de la Etapa 1$+$Elementos de la Etapa 2$+$Aplicaciones.

\end{itemize}
\newpage
\textbf{Tareas a realizar:}

\begin{enumerate}
    \item \textbf{Aplicaciones}: describir aplicaciones prácticas del tema investigado en el campo de la física o la ingeniería.
\end{enumerate}






\section*{Etapa 4 (P4): Presentaci\'on (3 \%)}

\begin{itemize}
    \item[\checkmark] Fecha de entrega apoyo audiovisual: martes 10 de junio 2025 (SEMANA 16).
    \item[\checkmark] Formato de entrega: documento en formato .pdf a través del módulo de evaluaciones en el \td.
    \item[\checkmark] Fecha de exposici\'on pública: martes 10 \'o jueves 13 de junio 2025 (SEMANA 16).
    \item[\checkmark] Duraci\'on:  m\'aximo 20 minutos por equipo.
\end{itemize}

\textbf{Tareas a realizar:}
\begin{enumerate}
    \item Preparar una exposici\'on pública del trabajo realizado, presentando los \textbf{principales conceptos} y \textbf{al menos una aplicaci\'on} del tema investigado en la F\'isica y la Ingenier\'ia. 
\end{enumerate}

\textit{Se evaluará la claridad en la presentación (incluyendo el apoyo audiovisual), el manejo del tema y del tiempo, así como la capacidad para responder preguntas del público. Adem\'as, se considerar\'a la participaci\'on de todo el equipo en la presentaci\'on.}

